\documentclass[11pt,a4paper]{article}
\usepackage[margin=2.2cm]{geometry}
\usepackage[T1]{fontenc}
\usepackage{amsmath,amssymb,xcolor,enumitem,booktabs,graphicx,tabularx,array}
\usepackage[normalem]{ulem}
\setlength{\parskip}{0.6em}\setlength{\parindent}{0pt}
\definecolor{newc}{HTML}{1f6f43}
\newcommand{\new}[1]{\textcolor{newc}{#1}}
\newcommand{\old}[1]{\textcolor{gray}{\sout{#1}}}
\newcommand{\note}[1]{{\small\itshape [#1]}}
\newcommand{\lab}[1]{\textsf{\footnotesize\detokenize{#1}}}
\renewcommand{\ref}[1]{\lab{#1}}
\renewcommand{\eqref}[1]{(\lab{#1})}
\newcommand{\nref}[1]{\lab{#1}}
\newcommand{\citep}[1]{[\lab{#1}]}
\newcommand{\citet}[1]{\lab{#1}}
\newcommand{\pitem}[2]{\subsection*{P#1. #2}}
\newcommand{\kind}[1]{\textbf{#1.}}
\title{Proposed change, Problem Formulation\\ (Checklist sections 5 and 6)}
\date{24 November 2026, version 1}
\begin{document}\sloppy
\maketitle

\section*{How to read this}
This answers sections 5 and 6 of \texttt{Checklist.md}, checked against the
current \texttt{3.ProblemFormulation.tex}, GAPS.md and the Method. Black text
stays, \new{green} is new and \old{struck} text goes. \note{Notes in
brackets} say why. Only the paragraphs that change are shown, in chapter
order. Nothing below is in the thesis yet.

Each item is marked \textbf{Fix} (a plain mistake or repeat) or
\textbf{Choice} (yours to decide, my recommendation is the green text).
Answer by number, for example ``P3 yes, P14 option b, P17 no''.

\section*{What the checklist asked that is already done or should not happen}
\begin{itemize}[nosep]
  \item \textbf{5.2 The environment.} Done. Generator details are not in
    the chapter, and ``a position is where an agent stands'' is said once.
    Only the repeated pointers to the Background are left (P1).
  \item \textbf{5.3 Collisions.} Done. The chapter states the two rules and
    one sentence on enforcement. Who yields, priority and chains of waiting
    are in the Method since [A126].
  \item \textbf{5.4 The three rules.} Each is one bullet. Scores, weights and
    the formula of $\hat\omega_t$ are in the Method.
  \item \textbf{5.5 Figures.} \lab{fig:storageexample} already sits beside the
    rules. \lab{fig:taskroute} stays. It holds the tea task $\tau_1$, which the
    life of a task, the stock condition, the observation at $t=12$ and
    \lab{fig:loop} all use. \lab{fig:storageupdate} stays, it is the only
    place a lost pallet is shown, and the checklist itself says lost pallets
    matter.
  \item \textbf{5.6 and 5.8 Pickup rule and controllers to the Method.} Not
    proposed. [A123] keeps the pickup rule in the Problem Formulation and out
    of the Method, and [A125] cut ``Choosing together'' to what the Problem
    Formulation does not say. Moving them back would undo both.
  \item \textbf{5.7 The subtraction $\zeta_1 = 14-5 = 9$.} Kept. It is one
    line, and your style asks for a worked number after a formula.
  \item \textbf{5.12 Count examples} (900 and 600 tasks, 800 orders and 10
    picks). Kept. Each makes one argument of the keep-up check concrete.
\end{itemize}

\section*{Summary}
\begin{tabular}{@{}llp{0.62\linewidth}@{}}
\toprule
Item & Kind & What \\
\midrule
P1 & Fix & Background named in full once per section, not five times \\
P2 & Choice & Pilot numbers out of the chapter, DSLAP named first \\
P3 & Choice & Capacity example in words, matrix product dropped \\
P4 & Choice & Checklist 6.3, a storage rule places stock \emph{and} ranks shelves \\
P5 & Choice & Idle agents move to ``A task's life'', busy and free said once \\
P6 & Fix & Lost text before \lab{eq:centralisedpolicy} \\
P7 & Fix & Forward pointers into the Method cut (distance label, crowding weight) \\
P8 & Choice & Crowding example on \lab{fig:observation} in one sentence \\
P9 & Fix & Subsection ``Putting the pieces together'' dropped, equation kept \\
P10 & Fix & Learning problem shorter, checklist 6.1 and 6.2 resolved \\
P11 & Choice & $\Delta$ paragraph trimmed to the trade-off \\
P12 & Fix & Lost text in the stall paragraph \\
P13 & Fix & Full stop after \lab{eq:split} \\
P14 & Choice & Gain test, \lab{eq:relgain} and the reporting line move to the Experiments, factor paragraph cut \\
P15 & Fix & Scope says the boundary only, ``If RQ4 is not reached'' cut \\
P16 & Choice & Delete about 200 lines of commented-out old text \\
\bottomrule
\end{tabular}

% ================================================================
\section*{The environment and Agents and collisions}

\pitem{1}{Back pointers to the Background (Fix)}
The chapter says ``described in Section~\ref{sec:bg:mapf}, Multi-Agent Path
Finding'' five times. The opening keeps both full pointers. Later mentions
lose them.

Every generated warehouse is well-formed in the sense of
\citet{ma2017lifelong_mapf}\old{, described in Section~\ref{sec:bg:mapf},
Multi-Agent Path Finding}.

This is the classical action model of multi-agent path finding
\citep{stern2019mapf}\old{, described in Section~\ref{sec:bg:mapf}, Multi-Agent
Path Finding}.

These are the vertex and swap conflicts \old{of Section~\ref{sec:bg:mapf}}
\new{described in the Background}.

This is the task of multi-agent pickup and delivery of
\citet{ma2017lifelong_mapf}\old{, described in Section~\ref{sec:bg:mapf},
Multi-Agent Path Finding}.

This is the dedicated storage \old{of Section~\ref{sec:bg:storage}}
\new{described in the Background}.

\note{The storage pointer in P2 moves with its sentence. Words, not
``Section III-B'', as you asked on 18 November.}

% ================================================================
\section*{Where the items are}

\pitem{2}{The storage rule paragraph (Choice)}
\textbf{The storage rule.}
A storage rule decides where arriving stock goes\old{, and nothing else}.
\new{This is the dynamic storage location assignment problem
\citep{WaubertdePuiseau2022DynamicSL}, described under Storage assignment in
the Background.}
Stock already on a shelf stays where it is until it is picked. This was
chosen over also moving stored stock, because every arriving pallet is a new
choice of shelf that needs no extra trip. Moving stored stock would need
trips of its own, or would have to assume they cost nothing. \old{A pilot
with Token Passing supports the choice. Demand storage that placed arriving
stock cut service time by about $43\%$ against fixed storage, and moving up
to a fixed number of stored units per period by $25\%$
(Appendix~\ref{app:cal:arrival}).} \old{This is the dynamic storage location
assignment problem \citep{WaubertdePuiseau2022DynamicSL}, described in
Section~\ref{sec:bg:storage}, Storage assignment.} Write $F$ for a storage
rule. The thesis compares three.

\note{Three things. (1) The problem is named first, then the choice. It came
after the pilot before. (2) ``and nothing else'' goes, see P4. (3) The pilot
numbers are a run result in the problem chapter, and the appendix pointer
points forward. [A123] moved them here today, so this undoes part of it.
Appendix G keeps the numbers. \textbf{Option b}: keep the pilot sentence but
drop the two percentages.}

\pitem{3}{The capacity example (Choice)}
\textbf{Example.}
Table~\ref{tab:storagematrix} shows a small storage state. Tea and mugs
take one unit of room each, and coffee takes two. Shelves $A$, $B$ and $C$
have room for 20, 15 and 10. \new{Shelf $A$ holds 12 tea and 4 coffee,
which take $12 + 8 = 20$, so it is exactly full. Shelf $B$ uses 12 of its 15
and $C$ uses 7 of its 10.} \old{Shelf $A$ is then exactly full, and $B$ and
$C$ have room left:}
\old{$(1\;2\;1)\cdot x = (20\;12\;7) \le (20\;15\;10)$} \note{the displayed matrix product}
Room also ties SKUs together. Tea and coffee share the room on shelf $A$,
so more tea there leaves less room for coffee.

\note{The checklist wanted the table and example in an appendix. Table I is
the base of your worked examples and of P4's storage update, so it stays.
Only the matrix product goes. The sum on one shelf shows the condition as
well and reads in words.}

\pitem{4}{What a storage rule is, checklist 6.3 (Choice)}
The text says a rule decides where arriving stock goes ``and nothing else''.
But congestion-aware storage also ranks the shelves an order line takes
from, and $F$ in \lab{eq:storageupdate} returns only a shelf per pallet. So
the decision problem \lab{eq:decisionproblem}, which ranges over
$\mathcal{F}$, does not cover the ranking that is compared.

\textbf{Option a (recommended).} A storage rule has two parts.

In \emph{The storage update}:
A storage rule is any $F$ that meets these three conditions, and
$\mathcal{F}$ is the set of them. \new{\note{unchanged for now, completed in
Orders and tasks below}}

In \emph{Under congestion-aware storage}, after ``its shelves are then listed
by $d_G(v_i,g)+\hat\omega_t(v_i)$ instead'':

\new{So a storage rule does two things. It places arriving stock at a
storage update, and it ranks the shelves a line takes from. Fixed and
demand storage rank by distance alone. From here on $F$ means both parts,
and $\mathcal{F}$ holds every rule that places stock under the three
conditions of a storage update and ranks the shelves for every line.}

\note{The notation table row for $\mathcal{F}$ then reads ``storage rules,
placement and shelf ranking''. The table of rules already has the column
``where a line takes stock from''. Nothing in the code changes.}

\textbf{Option b.} Keep $F$ as placement only, and say once that the
compared condition is a storage rule together with its shelf ranking. Then
\lab{eq:decisionproblem} ranges over those pairs and not over $\mathcal{F}$
alone.

% ================================================================
\section*{A task's life}

\pitem{5}{Busy and free said once, the idle agent moves here (Choice)}
\old{Only a free agent can get a new task. An agent is free when it has no
active task. Agent $a_2$ is busy from the moment it gets $\tau_1$, also on
its way to the shelf, and is free again once $\tau_1$ is delivered.}
\new{Only a free agent, one with no active task, can get a new task. So
agent $a_2$ is free again once $\tau_1$ is delivered.}

\new{A free agent stays where it is. When it stands in another agent's way,
it moves to the nearest free endpoint, until it gets there or gets a task.
In the way means on a shelf, workstation or dock of an unfinished task, or
on a vertex another agent was stopped from moving onto. Without this, an
idle agent would stay where its last task ended, often on a dock or a
workstation, and block the next agent that needs it.}

Then in \emph{Direction to the target}:

\old{An agent with no task also has a target. Usually it is the vertex it
stands on, so it stays put. When it stands in another agent's way, its
target becomes the nearest free endpoint until it gets there or gets a task.
In the way means on a shelf, workstation or dock of an unfinished task, or
on a vertex another agent was stopped from moving onto. Without this, an
idle agent would stay where its last task ended, often on a dock or a
workstation, and block the next agent that needs it.}
\new{An agent with no task has a target too, the vertex it stands on or,
when it is in the way, the nearest free endpoint, as in A task's life.}

\note{The first part is the repeat the checklist named in 5.7: ``busy from
assignment'' is already said two paragraphs up. The second part is what an
idle agent does, which is behaviour of the warehouse, not a feature of the
observation. [A97] put it in the Problem Formulation on purpose, and it
stays there, one section earlier. The checklist wanted it in the Method,
which would leave the model without it.}

% ================================================================
\section*{Who decides}

\pitem{6}{Lost text before the centralised policy (Fix)}
Formally, write $u_{\prec i}(t)$ for the moves already chosen by the agents
before $a_i$ in the priority order, where $\prec$ reads ``comes before''.
\new{The \emph{centralised policy} is}

\note{In the source the words follow a \texttt{\%~[A119]} tag on the same
line, so LaTeX drops them. The PDF reads ``comes before''. and then the
equation. The tag moves to the end of the line.}

% ================================================================
\section*{What a decentralised agent sees}

\pitem{7}{Pointers into the Method (Fix)}
In \emph{Direction to the target}:

The distance label is the shortest-path distance feature of
\citet{vanKnippenberg2021}. Because it is a distance and not just a
``closer or further'' flag, a policy can take a small step back for a bigger
gain later. It changes whenever the target changes. \old{The policy reads it
in two forms, as the difference from the label of the agent's own vertex
and as the label divided by a fixed number, so that it stays small
(Section~\ref{sec:method:model}).}

At the end of \emph{How crowded the neighbourhood is}:

\old{How strongly a learned policy reacts to $\delta_i(t)$ is set by its
weight in the reward. The weight is the same for every controller and is
not varied between runs (Section~\ref{sec:method:rl}).}

In \emph{The learning problem}:

Moves that lead nowhere are switched off before the policy chooses\old{
(Section~\ref{sec:method:model})}.

\note{The chapter points back only (your rule of 18 November). The two forms
of the label were moved to the Method in [A124], and the crowding weight is
in the reward table of [A127]. The pointer ``reward left open here'' in
\lab{tab:posg} stays, since it marks something left open on purpose.}

\pitem{8}{One crowding example on the observation figure (Choice)}
\old{For example, the field of view of $a_1$ in
Figure~\ref{fig:observation} covers $5$ other vertices besides its own, and
$2$ of them hold another agent ($a_2$ and $a_3$), so $\delta_1(t) = 2/5 =
0.4$. Two fifths of the nearby space is taken.}
\new{For agent $a_1$ in Figure~\ref{fig:observation}, two of the five other
vertices in view hold an agent, so $\delta_1(t) = 0.4$.}

\note{There are three crowding examples, this one, \lab{fig:crowding} and the
mean crowding in How a run is scored. \lab{fig:crowding} stays, because the
loop section and the mean crowding reuse it. This one becomes one sentence.}

\pitem{9}{``Putting the pieces together'' (Fix)}
\old{\textbf{Putting the pieces together} \note{subsection heading}}

The equation moves to the end of \emph{How crowded the neighbourhood is},
before Messages:

\new{The three pieces together form the \emph{decentralised observation} of
agent $a_i$,}
\[ o_i(t) = \bigl( \mathcal{G}^{(d_{\mathrm{obs}})}_i(t),\; \eta_i(\cdot, t),\; \delta_i(t) \bigr). \]

\note{A subsection of one equation. Messages then reads naturally, since it
says ``$o_i(t)$ keeps its three pieces''.}

% ================================================================
\section*{The learning problem}

\pitem{10}{Shorter, and the two contradictions resolved (Fix)}
\textbf{Inventory.}
\old{Every part of the game is something the reader has already met. The
agents are the fleet. The state $s_t$ is the one of
Section~\ref{sec:pf:controllers}, Who decides. It also holds the priority
order that conflict resolution uses. The actions are wait and move. The
transition rule turns a state and the agents' actions into the next state.
The observation is what each agent sees. The reward scores one step for one
agent, and the discount sets how much a later reward counts against one
now.} \new{The state $s_t$ is the one of Who decides, together with the
priority order that conflict resolution uses.} Formally, the learning problem
is \dots

\note{The equation, the sentence that names each symbol and
\lab{tab:posg} already list the parts. Only the priority order was new, so
it stays.}

\textbf{Policy and kernel.}
\old{The policy and the transition kernel answer two different questions.
The policy answers what each agent tries to do. From what the agent observes
it gives a probability to each legal action, for example move north with
probability $0.7$ and wait with probability $0.3$, and one action is drawn.
The drawn actions of all agents, side by side, form the joint action $u_t$,
the whole fleet's decision for this timestep. Under the centralised
controller the agents draw their actions one after another, each seeing the
earlier draws, and the result is the same kind of joint action. The
section-based controller does the same inside each zone.}
\new{The policy says what each agent tries to do. From what the agent sees it
gives a probability to each legal action, and one action is drawn. The drawn
actions of all agents form the joint action $u_t$.}

\old{The transition kernel answers what then happens to the warehouse. It
takes the current state and the joint action and applies everything that is
not a choice of any agent.} \new{The transition rule $P$ says what then
happens to the warehouse. It applies everything that is not an agent's
choice.} Conflicts are resolved, the shared assignment rule pairs free agents
with waiting tasks, the storage update is applied, and newly arrived orders
are added. The result is the next state $s_{t+1}$. \old{The kernel is
therefore not a state but the rule that produces one.} Its randomness is the
arrival of orders and the dock drawn for each shipment and put-away. Moves
have no randomness, so the same state and joint action always move the
agents in the same way. Formally, $P(s' \mid s_t, u_t)$ is the probability
that the next state is $s'$. The assignment rule and conflict resolution are
the same for every controller, so they are part of $P$, not of any policy.
\new{The learned policy $\pi_\theta$ is therefore only the routing rule
$\pi^{\mathrm{route}}$ of a controller.}

\note{Checklist 6.2. Who decides defines a controller as the pair of
assignment and routing, and this section puts assignment in $P$. The new
sentence ties the two. The pair stays, because Token Passing assigns its own
tasks. The sequential draw is already told in Who decides.}

\textbf{What changes between the controllers.}
\old{Only the observation changes between the controllers. Everything else in
the game is the same. A centralised agent sees the whole state, a
section-based agent sees its zone, and a decentralised agent sees its own
view.} \new{The controllers share the rest of the game. As Who decides set
out, they differ in three ways. They differ in what an agent sees, in
whether it knows the moves chosen before it, and in whose reward judges its
action. What an agent sees is the observation function.} Formally, the
\emph{observation function} is \dots

\note{Checklist 6.1. ``Only the observation changes'' contradicted Who
decides and the paragraph ``Sharing one objective'' further down, which
says the credit differs. The three views are in the equation, so the
sentence that lists them goes. ``The three views fit inside each other''
stays, it is the only place that says so.}

% ================================================================
\section*{How the two layers close the loop}

\pitem{11}{The storage period (Choice)}
\old{Separating the two clocks has one consequence that constrains the
study. Between two storage updates the routing traffic settles into a steady
pattern, and the storage rule only sees its average per step over the period
just ended. If $\Delta$ is too small, the traffic has not settled before the
next update, so $F$ reacts to noise, and the shelves the agents pick from
keep changing while they are still learning. If $\Delta$ is too large, the
numbers $\hat\rho$ and $\hat\omega$ are out of date by the time they are
used. The choice of $\Delta$ is therefore a parameter of the study, set by
a sweep (Table~\ref{tab:storageparams}), not a constant fixed in advance.}
\new{The two clocks make the storage period $\Delta$ a parameter of the
study. The storage rule sees only the average traffic of the period just
ended. If $\Delta$ is too short, the traffic has not settled before the next
update, so $F$ reacts to noise. If $\Delta$ is too long, $\hat\rho$ and
$\hat\omega$ are out of date by the time they are used.}

\note{The trade-off is part of the problem and stays. ``While they are still
learning'' is about training, and the sweep table points forward, so both
go. The Method's storage section already says the storage period is set in
Appendix G. \textbf{Option b}: move the whole paragraph to the end of the
Method's storage section and keep only the first new sentence here.}

% ================================================================
\section*{How a run is scored}

\pitem{12}{Lost text in the stall paragraph (Fix)}
A light load leaves longer quiet stretches, so the length grows with the
time between tasks. \new{It looks at the whole fleet and not at one} agent,
because one agent often waits by plan, as under Token Passing, while a
healthy fleet finishes tasks all the time.

\note{Same cause as P6, a \texttt{\%~[A118]} tag mid-line. The PDF reads
``\dots time between tasks. agent, because one agent \dots''.}

\pitem{13}{Full stop after the split (Fix)}
\lab{eq:split} ends with a comma and the next sentence starts ``The split is
taken over the picks''. The comma becomes a full stop.

\pitem{14}{The gain test to the Experiments (Choice)}
The checklist (5.12) asks to keep the objective here and move the evaluation
procedure. Most of it already is in the Experiments. ``Factors'' and
\lab{tab:expfactors} give the factors, ``Mapping questions to comparisons''
gives the RQs, $\Lambda_T/m$ and the percentage, and ``Scoring a condition''
gives the paired $t$-test. Three passages here repeat them.

\textbf{Stays in the Problem Formulation.} Keeping up \lab{eq:throughput},
the order service time \lab{eq:meanservice}, the split, crowding, the
reported measures, \lab{tab:measures} and the decision problem
\lab{eq:decisionproblem}. That is what the problem asks and what it costs.

At the end of the decision problem:

The thesis does not search for the best pair. It compares three rules under
three controllers\old{, and the gain test below asks whether a rule beats
fixed storage under one controller}\new{, each rule against fixed storage
under the same controller}.

Before \lab{tab:measures}:

\old{Every quantity is reported as a mean over evaluation seeds with a 95\%
confidence interval, as \citet{krnjaic2022warehouse_marl} report their pick
rates.}

The paragraphs \emph{The research question, formally} (\lab{eq:rqformal} and
\lab{eq:relgain} with their text) move to the Experiments, as a paragraph
\emph{The gain test} right before ``Mapping questions to comparisons''. The
text is unchanged, apart from the first sentence and the reporting line,
which joins it:

\new{\textbf{The gain test.} The comparison is paired. Under one seed, fixed
storage and the rule $F$ receive the same customer orders.} The tasks can
still differ, because where a line is picked and how many tasks it splits
into depend on where the stock is. \dots\ \note{rest as now, with both
equations and their labels} \new{Every quantity is reported as a mean over
evaluation seeds with a 95\% confidence interval, as
\citet{krnjaic2022warehouse_marl} report their pick rates.}

The last paragraph of the section goes:

\old{The study then asks how the gain depends on three factors. The first is
the controller, how many agents one decision of $\pi$ covers. The gain is
read together with its split, because a fleet that runs close to its limit
turns the same saving in travel into a larger cut in the wait before
assignment. The second is whether the storage rule is congestion-aware,
congestion-aware storage against demand storage. The third is the load, set
by the fleet size $m$ and the expected number of order lines released per
timestep $\lambda_{\mathrm{task}}$. On the same sweep, throughput per agent,
$\Lambda_T / m$, shows whether each agent finishes fewer tasks as the fleet
grows, and whether that loss is faster when one decision covers fewer
agents. For RQ4, the decentralised controller is also run with
communication over \nref{eq:commgraph} and without it. Only that controller
has messages, so this comparison sits outside the main factorial design.}

\note{After [A129] the factors are told in Aim and scope, and again in the
Experiments' factor table and mapping paragraph. This was the third copy.
One sentence in it is not said elsewhere: a fleet close to its limit turns a
saving in travel into a larger cut in waiting. Aim and scope says nearly
the same (``Near that limit a small saving in travel turns into a large
saving in waiting''), so nothing is lost. The labels \lab{eq:rqformal} and
\lab{eq:relgain} keep their names, so the Experiments' references still
work, and the Problem Formulation no longer points to them. \textbf{Option
b}: leave the gain test here and only cut the factor paragraph and the
reporting line.}

% ================================================================
\section*{Scope and assumptions}

\pitem{15}{The boundary, not how to extend it (Fix)}
Tasks carry no deadline. A late delivery only lengthens service time, and no
single task can be said to have failed. This is narrower than
time-constrained formulations such as \citet{vanKnippenberg2021}, where a
missed deadline is infeasible or penalised. \old{Adding deadlines would put a
deadline in every task, and either forbid a late task or count late tasks
next to the result.}

Assignment is also permanent once made, so an agent serving a task is not
re-tasked before it finishes. \citet{ma2017lifelong_mapf} allow re-tasking
until the agent reaches where the task picks up. \old{Adding it would let the
assignment rule change the agent of a task whose agent has not yet reached
the shelf. Either change alters what counts as a solution, so it would
change this chapter and not the method.} \new{Either change would alter what
counts as a solution.}

Communication is built in and can be switched on or off\old{, but RQ1 to RQ3
do not depend on it. If RQ4 (Section~\ref{sec:pd-rq}) is not reached, the
results say nothing about communication.}\new{. Only RQ4 uses it.}

\note{Checklist 5.13 and the open item of [A129]. RQ4 is run, so the
conditional goes. How to add a deadline or re-tasking belongs to future
work, if anywhere.}

% ================================================================
\section*{Housekeeping}

\pitem{16}{Old text in comments (Choice)}
About 200 lines of commented-out old text sit in the chapter, lines 132 to
184 (old collisions and well-formedness) and 409 to 575 (old storage
section with $D+\beta K$). They do not appear in the PDF. The old text is in
git. Delete them?

% ================================================================
\section*{When approved}
Applied under a new tag (next free is [A132]), with a GAPS.md entry. If P4a
is chosen, the notation table row of $\mathcal{F}$ changes in both notation
tables. If P14 is chosen, the moved paragraph is checked against the
Experiments' own text so the gain is not explained twice there. Every tag
goes at the end of a line. The chapter is rebuilt and read back with
pdftotext.

\end{document}
